% !TeX spellcheck = en_US
\documentclass[acmsmall,natbib,review,screen,anonymous,10pt]{acmart}
% TODO: FINALIZE
\settopmatter{printfolios=true,printccs=false,printacmref=false}

%\newcommand{\commentaire}{}
\newcommand{\MA}[1]{\ifdefined\commentaire{{{\color{green!60!black!55} 
\texttt{[#1]}}}}\fi}
\newcommand{\GM}[1]{\ifdefined\commentaire{{{\color[HTML]{5100FF} 
\texttt{[#1]}}}}\fi}
\newcommand{\GV}[1]{\ifdefined\commentaire{{{\color[HTML]{5100FF} 
\texttt{\small[#1]}}}}\fi}
\newcommand{\gb}[1]{\ifdefined\commentaire{{{\color{green!60!black!55} 
\texttt{[#1]}}}}\fi}
\newcommand{\bg}[1]{\ifdefined\commentaire{{{\color{green!60!black!55} 
\texttt{[#1]}}}}\fi}
\newcommand{\blue}[1]{\color[HTML]{0008FF} #1}


\newcommand*{\EHL}{\texttt{eHL}\xspace} %MA: abbreviation for expectation logic
\newcommand*{\HL}{\texttt{HL}\xspace}
\newcommand*{\PRHL}{\texttt{pRHL}\xspace}
\newcommand*{\ec}{\texttt{EasyCrypt}\xspace}
\newcommand{\easycrypt}{\ec}

\acmJournal{PACMPL}
\acmVolume{1}
\acmNumber{CONF} % CONF = POPL or ICFP or OOPSLA
\acmArticle{1}
\acmYear{2018}
\acmMonth{1}
\acmDOI{} % \acmDOI{10.1145/nnnnnnn.nnnnnnn}
\startPage{1}
\setcopyright{none}

\bibliographystyle{ACM-Reference-Format}
\citestyle{acmauthoryear}

\usepackage[utf8]{inputenc}
\usepackage{amsmath}
\usepackage{hyperref}
\usepackage{booktabs}
\usepackage{caption}
\usepackage{wrapfig}
\usepackage{subcaption}
\usepackage{multicol}
\usepackage{xspace}
\usepackage{xfrac}
\usepackage{tikz}
\usetikzlibrary{matrix,backgrounds,arrows.meta}
\usepackage{cleveref}
\usepackage{sbnf}
\usepackage{paper}

\makeatletter
\let\@msm@th@eqref\eqref
\renewcommand{\eqref}[1]{%
  \begingroup
  \leavevmode
  \color{violet}%
  \hypersetup{linkbordercolor=[named]{violet}}%
  \@msm@th@eqref{#1}%
  \endgroup
}
\makeatother

\newcommand{\eqtag}[1]{\tag{{\color{violet}{#1}}}}

\captionsetup{subrefformat=parens}

\begin{document}

\title{Lower bound reasoning with \textsf{eHL}}

\maketitle

\begin{figure}
    \[
      \begin{array}{cc}
        \multicolumn{2}{c}{
        \Infer[hoare][skip]
        {\HT{\pred}{\ex}{\SKIP}{\pred}{\ex}}
        {}}
        \\[3mm]
        \multicolumn{2}{c}{
        \Infer[hoare][assign]
        {\HT{\pred[\vx \mapped \expr]}{\ex[\vx \mapped \expr]}{\vx <- \expr}{\pred}{\ex}}
        {}}
        \\[3mm]
        \multicolumn{2}{c}{
        \Infer[hoare][sample]
        {\HT{\forall \val \in \supp(\sexpr),\ \pred[\vx \mapped \val]}{\E{\sexpr}{\lambda \val.\, \ex[\vx \mapped \val]}}{\vx <* \sexpr}{{\pred}}{\ex}}
        {}}
        \\[3mm]
        \multicolumn{2}{c}{
        \Infer[hoare][seq]
        {\HT{\pred}{\ex}{\cmd;\cmdtwo}{\predtwo}{\extwo}}
        {\HT{\pred}{\ex}{\cmd}{\predthree}{\exthree} & \HT{\predthree}{\exthree}{\cmdtwo}{\predtwo}{\extwo}}}
        \\[3mm]
        \multicolumn{2}{c}{
        \Infer[hoare][if]
        {\HT{\pred}{\ex}{\ITE{\bexpr}{\cmd}{\cmdtwo}}{\predtwo}{\extwo}}
        {\HT{\bexpr\land\pred}{\ex}{\cmd}{\predtwo}{\extwo}& \HT{\neg\bexpr\land\pred}{\ex}{\cmdtwo}{\predtwo}{\extwo}}}
        \\[3mm]
        \multicolumn{2}{c}{
        \Infer[hoare][while]
        {\HT{\pred}{\lub\ex_{k}}{\WHILE\bexpr\DO\cmd}{\neg\bexpr\land\pred}{\extwo}}
        {\bexpr \land \pred \valid \ex_{0} \HTrel 0 
        & \forall k.\, \HT{\bexpr\land\pred}{\ex_{k+1}}{\cmd}{\pred}{\ex_{k}}
        & \neg\bexpr \land \pred \valid \ex_{k} \HTrel g}
        }
      \end{array}
    \]
  \caption{\EHL for lower and upper-bounds. Here ${\rel} \in \{{\leq},{\geq}\}$.}\label{fig:hoare:l-u}
\end{figure}

\begin{figure}
    \[
      \Infer[hoare][while]
      {\HT[lower]{\pred}{\ex}{\WHILE\bexpr\DO\cmd}{\neg\bexpr\land\pred}{\ex}}
      {
        \HT[lower]{\bexpr\land\pred}{\ex}{\cmd}{\pred}{\ex}
        & \bexpr \land \pred \valid r_0 \geq \ex
        & \forall i \geq 0.\, \HT[upper]{\bexpr\land\pred}{r_{i+1}}{\cmd}{\top}{[\pred]\cdot r_i}
        & \bexpr \land \pred \valid \lim_{i \to \infty} r_i = 0
      }
    \]
  \caption{The \textsc{While} rule for \textsf{eHL} lower bounds. The rule is
  specific to the lower-bound judgment: Premise~1
  is a lower-bound triple on the body, while Premise~3 is an
  \emph{upper}-bound triple used purely to bound the in-loop mass. The sequence $\{r_i\}_{i \geq 0}$
  consists of non-negative expectations $r_i : \mathrm{State} \to
  \mathbb{R}^+$.}\label{fig:hoare:while-lb}
\end{figure}


\paragraph{Notation.}
For a predicate $\pred$ we use two indicator-style operators, one tailored to
each direction of reasoning:
\[
  [\pred] \;=\;
  \begin{cases} 1 & \text{if } \pred \\ 0 & \text{otherwise} \end{cases}
  \qquad\text{(used for lower bounds)}
  \\[3mm]
\]
\[
  \langle \pred \rangle \;=\;
  \begin{cases} 1 & \text{if } \pred \\ \infty & \text{otherwise} \end{cases}
  \qquad\text{(used for upper bounds)}
  \\[3mm]
\]
Given a (state-indexed) expectation $r$, the expressions $[\pred] \cdot r$ and
$\langle \pred \rangle \cdot r$ denote the corresponding pointwise products:
$[\pred] \cdot r$ agrees with $r$ on states satisfying $\pred$ and is zero
elsewhere, while $\langle \pred \rangle \cdot r$ agrees with $r$ on states
satisfying $\pred$ and is $\infty$ elsewhere.

\paragraph{The witness sequence $\{r_i\}_{i \in \mathbb{N}}$.}
The rule of \Cref{fig:hoare:while-lb} uses a sequence $\{r_i\}_{i \geq 0}$ of
non-negative expectations as a termination witness: intuitively, $r_i$ is a
(user-chosen) upper bound on how much of the post-expectation mass of $\ex$ is
still trapped inside the loop after $i$ body executions. Premise~2 ties $r_i$
to $\ex$ at $i=0$ by requiring $r_0 \geq \ex$; Premise~3 propagates the bound
across one iteration of the body; and Premise~4 forces the bound --- and thus
the trapped mass --- to vanish in the limit.


\end{document}
